Распространение термоупругих волн в геологических средах описывается системой связанных уравнений термоупругости, объединяющих уравнения движения и уравнения теплопереноса. В классической линейной модели Куна–Томпсона используется система:
\begin{equation}
\rho \frac{\partial^2 u_i}{\partial t^2} = \frac{\partial \sigma_{ij}}{\partial x_j},
\label{eq:motion}
\end{equation}

где $u_i$ — компоненты вектора смещения, $\rho$ — плотность среды, $\sigma_{ij}$ — компоненты тензора напряжений.

Термоупругая связь между деформациями и напряжениями задаётся:

\begin{equation}
\sigma_{ij} = \lambda \delta_{ij} \epsilon_{kk} + 2\mu \epsilon_{ij} - \gamma \delta_{ij} T,
\label{eq:stress}
\end{equation}

где $\lambda, \mu$ — параметры Ламе, $\epsilon_{ij}$ — тензор деформаций, $\gamma$ — коэффициент термоупругой связи, $T$ — изменение температуры.

Уравнение теплопереноса с учетом термомеханической генерации записывается как:

\begin{equation}
k \nabla^2 T - C \frac{\partial T}{\partial t} = \gamma T_0 \frac{\partial \epsilon_{kk}}{\partial t},
\label{eq:heat}
\end{equation}

где $k$ — коэффициент теплопроводности, $C$ — теплоёмкость, $T_0$ — средняя температура.

Данная система уравнений является базовой для численного моделирования и построения обучающих выборок для алгоритмов машинного обучения.